\documentclass[10pt,a4paper]{amsart}
\usepackage[margin=19mm]{geometry}
\usepackage{amsmath,amssymb,mathtools}
\usepackage[scr=rsfs,cal=euler]{mathalfa}
\usepackage{xfrac}
\usepackage[unicode,hidelinks,bookmarks=false]{hyperref}
\newcommand{\ie}{i.e.\ }
\newcommand{\cf}{cf.\ }
\newcommand{\resp}{respectively\ }
\newcommand{\Subsection}{\subsection}
\providecommand{\nobreakdash}{\nobreak\hbox{-}}
\setlength{\emergencystretch}{2em}
\multlinegap0pt
\usepackage{xcolor}
\usepackage{amssymb}
\usepackage{tikz}
\usepackage{algorithm}\usepackage{algpseudocode}
\usepackage[shortlabels]{enumitem}
\newtheorem{lemma}{Lemma}
\newcommand{\Ahat}{\widehat{A}}
\newcommand{\Astar}{A}
\newcommand{\Bhat}{\widehat{B}}
\newcommand{\Bstar}{B}
\newcommand{\Ebound}{\mathcal{E}_{\mathrm{bound}}}
\newcommand{\Els}{\mathcal{E}_{\mathrm{ls}}}
\newcommand{\Esafe}{\mathcal{E}_{\mathrm{safe}}}
\newcommand{\Fnorm}[1]{\left\lVert #1 \right\rVert_{\mathrm{F}}}
\newcommand{\Vfast}{V_{\mathrm{fast}}}
\newcommand{\Vfasthat}{\widehat{V}_{\mathrm{fast}}}
\newcommand{\Vslow}{V_{\mathrm{slow}}}
\newcommand{\Vslowhat}{\widehat{V}_{\mathrm{slow}}}
\newcommand{\cfast}{c_{\mathrm{fast}}}
\newcommand{\cslow}{c_{\mathrm{slow}}}
\newcommand{\sbase}{s_k^{\mathrm{base}}}
\newcommand{\tauk}{\tau_k}
\title{Quantized Online LQR}
\author{Barron Han, Victoria Kostina, Babak Hassibi}
\date{Source: arXiv:2604.11930v1 (CC BY 4.0)}
\begin{document}
\maketitle

\noindent\emph{Source and licence.} Quantized Online LQR. Version arXiv:2604.11930v1 (CC BY 4.0), distributed by arXiv under the Creative Commons Attribution 4.0 International licence, http://creativecommons.org/licenses/by/4.0/, as recorded in the arXiv licence metadata for this exact version and shown on the abstract page https://arxiv.org/abs/2604.11930v1. Full licence notice: \texttt{licenses/arxiv-2604.11930-CC-BY-4.0.txt}.\par\smallskip

\noindent\textit{Editorial scope.} Counted unit: the article's lemma lem:contraction (source0.tex lines 609--616). The article's complete original proof of it is delivered with it (source0.tex lines 1336--1355, one \texttt{proof} environment of 20 lines, which the article places away from the statement). No mathematics is written, repaired or re-derived: the statement and the proof are the article's own lines, in the article's own notation, and no retained line is substituted or re-flowed.  Retained uncounted as the source-local context the proof needs: lines 553--556; lines 593--596.  Nothing was omitted from any retained span, so the package carries no omission record.  No retained line contains a citation, so the wrapper carries no bibliography and no bibliography entry.  No retained line contains a figure environment, an image-inclusion command or drawing code, so no image is delivered and the package declares no asset dependency.  Editorial formatting only, in addition: an AMS \texttt{amsart} preamble that restates the article's own declarations and macros used by the retained lines, namely the environments \texttt{lemma} on separate counters, together with the standard packages those lines need.  The retained environments print in this document as Lemma 1 in order of appearance, whereas the article prints the same items with the numbering of its own full document.  The complete native source of the article as distributed in the arXiv e-print for this exact version is bundled unchanged as \texttt{sources/198370/source0.tex}.
    \begin{equation}
        \label{eq:ols-two-rate}
        \Fnorm{\Ahat_k - \Astar} + \Fnorm{\Bhat_k - \Bstar} \le \Vslow\,\tauk^{-1/4} + \Vfast\,\tauk^{-1/2},
    \end{equation}

\begin{align}
    \cslow &= \frac{1 + 2^{1/4}}{1 - \varrho\,2^{1/4}}\,\Vslowhat, \label{eq:cslow-impl} \\
    \cfast &= \frac{1+\sqrt{2}}{1 - \varrho\,\sqrt{2}}\,\Vfasthat. \label{eq:cfast-impl}
\end{align}

\begin{lemma}[Contraction of the Adaptive Multiplier]
    \label{lem:contraction}
    Let $\varrho < 1/\sqrt{2}$ and $\beta := \varrho\sqrt{2} < 1$. On the event $\Els \cap \Esafe \cap \Ebound$, for all epochs $k > k_{\mathrm{ls}}$ (where the asymptotic OLS bound \eqref{eq:ols-two-rate} holds), the multiplier satisfies
    \begin{equation}
        m_k < 1 + \beta(m_{k-1} - 1) + 1.
    \end{equation}
    Consequently, $\limsup_{k\to\infty} m_k \le \frac{2 - \beta}{1-\beta} = \mathcal{O}(1)$.
\end{lemma}

\begin{proof}[Proof of Lemma~\ref{lem:contraction}]
    For $k > k_{\mathrm{ls}}$, the OLS bounds hold. Let $\mathrm{OLS}_k := \Vslowhat\,\tau_k^{-1/4} + \Vfasthat\,\tau_k^{-1/2}$. By the triangle inequality centered at $\theta_\star$ and non-expansiveness of projection:
    \[
        \|\Delta_k\|_2 \le \|\bar\theta_k - \theta_\star\|_2 + \|\theta_\star - \bar\theta_{k-1}\|_2 + \|\bar\theta_{k-1} - \tilde\theta_{k-1}\|_2 \le \mathrm{OLS}_k + \mathrm{OLS}_{k-1} + \varrho s_{k-1}.
    \]
    By the design of $\cslow$ and $\cfast$ in \eqref{eq:cslow-impl}--\eqref{eq:cfast-impl}:
    \[
        \mathrm{OLS}_k + \mathrm{OLS}_{k-1} \le \sbase - \varrho s_{k-1}^{\mathrm{base}}.
    \]
    Indeed, using $\tau_{k-1}^{-1/4} = 2^{1/4}\tau_k^{-1/4}$ and $\tau_{k-1}^{-1/2} = \sqrt{2}\,\tau_k^{-1/2}$, the left-hand side equals $\Vslowhat(1+2^{1/4})\tau_k^{-1/4} + \Vfasthat(1+\sqrt{2})\tau_k^{-1/2}$, while the right-hand side expands to $\cslow(1-\varrho\,2^{1/4})\tau_k^{-1/4} + \cfast(1-\varrho\sqrt{2})\tau_k^{-1/2}$, and by the definitions \eqref{eq:cslow-impl}--\eqref{eq:cfast-impl} these are identical.
    Substituting $s_{k-1} = m_{k-1} s_{k-1}^{\mathrm{base}}$:
    \[
        \|\Delta_k\|_2 \le \sbase + \varrho s_{k-1}^{\mathrm{base}}(m_{k-1} - 1).
    \]
    Dividing by $\sbase$ and using $s_{k-1}^{\mathrm{base}}/\sbase \le \sqrt{2}$ (since the ratio is a weighted average of $2^{1/4}$ and $\sqrt{2}$ with positive weights $\cslow\tau_k^{-1/4}$ and $\cfast\tau_k^{-1/2}$, hence bounded by $\sqrt{2}$):
    \[
        \frac{\|\Delta_k\|_2}{\sbase} \le 1 + \varrho\sqrt{2}(m_{k-1}-1).
    \]
    Since $m_k = \lceil \|\Delta_k\|_2/\sbase \rceil$ and $\lceil x \rceil < x + 1$, the contraction follows. With $\beta = \varrho\sqrt{2} < 1$, the sequence converges exponentially to the fixed point $m^\star = (2-\beta)/(1-\beta) = \mathcal{O}(1)$.
\end{proof}

\end{document}
